\section{Statements, normalization, and the scope of the comparison}
\label{b-introduction}

The question addressed here is whether adding finitely many prescribed
positive sources can lower the ground energy of a cutoff QED model
by more than a constant times their number, after subtracting the
\emph{same model's interacting neutral vacuum}. This is distinct from
an absolute lower bound in the free-vacuum normalization. We prove a
relative theorem for a sharply specified momentum regulator, while
also proving that this regulator has a superextensive absolute
vacuum energy. Both conclusions concern the same Hamiltonian.

The model has finite fermionic one-particle dimension and finitely
many nonzero photon modes in a torus, but unbounded photon occupation.
Its fermionic Fock space permits arbitrary electron--positron pair
creation up to the regulator's finite capacity. A sea representation
on $\bigwedge\mathfrak h$ is unitarily equivalent to the electron and
positron Fock representation; it is not an unbounded-below
first-quantized Dirac prescription. The exact conserved quantity is
the electron excess, including the prescribed nuclear charge in the
neutrality condition. A mean-charge condition on a mixed quasi-free
state is distinguished from membership in one exact charge sector.

The fermion band is fixed independently of the electromagnetic
field. Local multiplication by a gauge phase generally does not
preserve it. Thus this is a Coulomb-gauge cutoff Hamiltonian with
an imposed global Gauss constraint, not a locally gauge-covariant
continuum construction. The harmonic transverse zero mode is
explicitly omitted; Coulomb gauge by itself does not remove that
degree of freedom. These choices are relevant to the counterexamples.

In the full free-sea Wick model the source-free Hamiltonian can be
rewritten as a positive sum of centered charge-density squares plus
a one-body Dirac operator and a fixed scalar. The one-body mass is
\[
 M_{\rm eff}(p)=m-\frac{\alpha}{2L^3}
 \sum_{\substack{r\in(2\pi/L)\Z^3,\ |r|\le\Lambda_f\\r\ne p}}
 \frac{4\pi}{|p-r|^2}\frac{m}{\sqrt{m^2+|r|^2}}.
\]
When $M_{\rm eff}\ge m_*I>0$, the centered fermion trace in a
real photon-coordinate history and arbitrary scalar auxiliary-field
kicks is positive. Its transfer eigenvalues lie outside an explicit
annulus. A canonical Fourier estimate then preserves the desired
charge sector rather than using a grand-canonical bound in its place.
The proofs below do not replace noncommuting density operators by
classical variables.

For distinct point nuclei, integral $1\le Z_i\le Z_{\max}$,
$Z_{\rm tot}\le2M_L$, and $L\ge\Lambda_f^{-1}$, the resulting
full-field inequality is
\begin{equation}
 E(R,Z)-E^{\vac}\ge m_*Z_{\rm tot}
 -\alpha\Lambda_f\left[
 20\pi(1+1/\pi)^3Z_{\max}
 +\frac{13}{4\pi^{3/2}}Z_{\max}^2\right]K.
 \label{b-headline}
\end{equation}
The proof is Theorem~\ref{pe-full-theorem} with the strengthening
of Theorem~\ref{thm:mass-gap-charge}. A single nucleus instead obeys
$E_Z-E^{\vac}\ge m_*Z-13\alpha\Lambda_fZ^2/\pi^2$ for every
admissible box. Nuclear deltas are preserved in the source definition
and in the exact nuclear repulsion. Compression of their coupling
to the fermion band regularizes the potential seen by the field.
The entire individual periodic nuclear diagonal is subtracted by a
fixed prescription. No state-dependent or configuration-dependent
subtraction is used.

One simple sufficient mass condition is
$13\alpha\Lambda_f<\pi^2m$. Section~\ref{sec:mass-gap-round2}
proves a larger all-volume window, including
$\Lambda_f/m<5168.21$ at $\alpha=1/137$, as an evaluation of an
explicit inequality. A continuum logarithmic coefficient is proved
with a mesh error, and a coarse-box example shows why that error
cannot be omitted uniformly. Density-product ordering removes this
particular mass-window restriction by changing a definite one-body
operator, not merely a scalar normalization.

An absolute extensive statement would instead read
\begin{equation}
 H_{\Lambda,L}(R,Z)\ge-C_\Lambda(K+L^3).
 \tag{W}\label{eq:W}
\end{equation}
Section~\ref{sec:counterexamples} disproves it for the bare momentum
regulator with massless photons in its specified free-vacuum normalization. A
fixed positive transverse photon-mass deformation restores such a
bound, with the squeezing constant retained. Static susceptibility
matching and its Ward-identity interpretation are treated separately
from this operator theorem. The subsequent mean-field sections prove
two further obstructions and carefully distinguish their reference
vacua from the unrestricted correlated vacuum used in
\eqref{b-headline}.

Notation in the variational sections uses $P_-=P$, $P_+=I-P$,
$\Gamma_{f,L}=\Gamma_{\Lambda_f,L}$, and $\mathfrak h_f=\mathfrak h$.
There, a one-particle symbol $Q=\gamma-P_-$ is a sea-relative density
matrix; elsewhere the conserved Fock charge is $Q=N-2M_L$.
Their meanings are stated where they are used. All bare parameters
are fixed unless a limit is explicitly declared.
